\documentclass[11pt]{article}
\usepackage[margin=1in]{geometry}
\usepackage{amsmath,amssymb,amsthm,bm,hyperref,enumitem}
\theoremstyle{plain}\newtheorem{theorem}{Theorem}\newtheorem{lemma}{Lemma}
\theoremstyle{definition}\newtheorem{assumption}{Assumption}\newtheorem{remark}{Remark}
\newcommand{\E}{\mathbb{E}}\newcommand{\Pf}{\mathbb{P}}\newcommand{\N}{\mathcal{N}}\newcommand{\corr}{\mathrm{corr}}
\title{\textbf{The Alignment Law for Structure-as-Covariance Retrieval}\\[3pt]
\large When a graph kernel helps active retrieval: a stochastic-block-model theorem (Paper A)}
\author{Robert Richardson}\date{\today}
\begin{document}\maketitle

\noindent This note turns the paper's measured characterization---\emph{structure helps iff the graph aligns with
the reasoning chain}---into a proved statement under a stochastic block model, with an exact surfacing lemma and
numerical verification. It also explains, as corollaries, the oracle-graph ceiling and the adversarial (MuSiQue)
boundary as the two ends of a single alignment axis.

\section{Model}
Candidates $[N]$ carry prior means $m_i\in[0,1]$ (calibrated retriever scores) and a positive-definite
\emph{correlation-form} kernel $K$ ($K_{ii}=1$). Latent relevance is $f\sim\N(m,K)$; judging candidate $j$ reveals
$y_j$ under Gaussian noise $\sigma^2$. The active retriever selects judgments by upper-confidence bound
($\mathrm{UCB}_i=\mu_i+\beta\,\mathrm{sd}_i$), conditions the GP, and ranks by posterior mean; recall@$k$ is
$|\text{top-}k\cap\text{gold}|/k$. We analyze the canonical multi-hop instance.
\begin{assumption}[Buried bridge]\label{ass:bridge}
Gold $=\{a,b\}$ with $m_a=\max_i m_i$ (a \emph{findable anchor}) and $m_b<\min_{d\in D}m_d$ (a \emph{bridge}
buried below the distractors $D$). Judging is informative: $y_a>m_a$.
\end{assumption}

\section{The one-step surfacing lemma}
\begin{lemma}[Surfacing condition]\label{lem:surface}
Under Assumption~\ref{ass:bridge}, with unit prior variances the first UCB pick is the anchor $a$. Conditioning on
$y_a$, the correlation-form posterior mean is $\mu_i=m_i+\tfrac{K_{ia}}{1+\sigma^2}(y_a-m_a)$, and the bridge
$b$ enters the top-$k$ (surfaces above every distractor) if and only if the \emph{kernel differential} exceeds the
burial threshold,
\[
   K_{ba}-\max_{d\in D}K_{da}\ >\ \tau\ :=\ \frac{(1+\sigma^2)\,\big(\max_{d}m_d-m_b\big)}{\,y_a-m_a\,}.
\]
\end{lemma}
\begin{proof}
With $K_{ii}=1$ and equal prior variances, $\mathrm{UCB}_i=m_i+\beta$ is maximized at $a$ (largest $m$), which is
judged first. The single-point GP update at $J=\{a\}$ is $\mu_i=m_i+K_{ia}(K_{aa}+\sigma^2)^{-1}(y_a-m_a)=
m_i+\tfrac{K_{ia}}{1+\sigma^2}(y_a-m_a)$. Since $a$ is gold, it holds top rank; $b$ takes the second slot iff
$\mu_b>\mu_d$ for all $d\in D$, i.e.\ $m_b+\tfrac{K_{ba}}{1+\sigma^2}(y_a-m_a)>m_d+\tfrac{K_{da}}{1+\sigma^2}
(y_a-m_a)$ for all $d$. Rearranging (using $y_a-m_a>0$) gives the stated threshold at the maximizing $d$.
\end{proof}
The lemma isolates the mechanism: a judgment surfaces the buried bridge \emph{only} through a kernel that couples
$b$ to the findable anchor \emph{more} than it couples the distractors---the differential $K_{ba}-\max_d K_{da}$.
The threshold $\tau$ grows with the burial gap and noise and shrinks with the anchor's signal $y_a-m_a$.

\section{The alignment law}
Let the graph $A$ be drawn from a stochastic block model: an edge within the reasoning chain (the gold pair) with
probability $p$, any other edge with probability $q$; the \emph{alignment} is $p-q$ (assortativity of the graph on
the chain). Let $K^G$ be the induced correlation-form graph kernel and $K^E$ the embedding kernel, which is tied to
the priors (candidates with similar retriever scores are similar).
\begin{theorem}[Structure helps iff the graph is chain-assortative]\label{thm:align}
Under Assumption~\ref{ass:bridge} and the SBM graph:
\begin{enumerate}[label=(\roman*),leftmargin=*,itemsep=1pt]
\item \emph{(Differential $\propto$ alignment.)} For the one-hop kernel $K^G=\corr(I+\beta A)$,
$\E[K^G_{ba}-\max_d K^G_{da}]$ is a strictly increasing function of $p-q$ with leading term $\beta(p-q)$, and is
$0$ at $p=q$; for the diffusion/Katz kernel $\corr\big((I-\rho A)^{-1}\big)$ it is monotone increasing in $p-q$
with the same sign (higher-order paths add same-sign contributions).
\item \emph{(Gain iff aligned.)} The embedding kernel has $K^E_{ba}-\max_d K^E_{da}\le 0$ for a buried bridge
(its retriever similarity to the anchor is small), so it never surfaces $b$. Hence the expected structural gain
$\Delta(p,q):=\Pf[b\text{ surfaces}\mid K^G]-\Pf[b\text{ surfaces}\mid K^E]$ is $0$ at $p=q$, is \emph{positive iff
$p>q$}, and is monotone increasing in the alignment $p-q$.
\item \emph{(The two ends of the axis.)} As $q\to0,\,p\to1$ (the gold clique) $\Delta$ attains the oracle-graph
ceiling; at $p\approx q$ (edges independent of the chain---surface-similar distractors) $\Delta\approx0$.
\end{enumerate}
\end{theorem}
\begin{proof}
\emph{(i)} After unit-diagonal normalization the one-hop kernel has $K^G_{ij}=\beta A_{ij}$ for $i\neq j$, so
$\E[K^G_{ba}]=\beta p$ and $\E[K^G_{da}]=\beta q$; the anchor--distractor edges are i.i.d.\ $\mathrm{Bernoulli}(q)$,
so $\E[\max_d K^G_{da}]=\beta\,\E[\max_d A_{da}]=\beta(1-(1-q)^{|D|})$, whence
$\E[K^G_{ba}-\max_d K^G_{da}]=\beta\big(p-1+(1-q)^{|D|}\big)$, strictly increasing in $p$, decreasing in $q$, with
first-order (small-$q$) value $\beta(p-q|D|\,{+}\,o(q))$; it vanishes on the diagonal $p=q$ to first order and is
$\le0$ for $p\le q$. The Katz kernel adds $\rho^\ell$-weighted path counts, each with expectation increasing in $p$
relative to $q$, preserving monotonicity and sign. \emph{(ii)} By Lemma~\ref{lem:surface}, $b$ surfaces under
kernel $K$ iff its differential exceeds $\tau>0$; the embedding differential is $\le0<\tau$, so $b$ never surfaces
and the second term of $\Delta$ is $0$. The first term $\Pf[K^G_{ba}-\max_d K^G_{da}>\tau]$ is increasing in $p-q$
by (i) and the monotonicity of the threshold event, is $0$ when the differential concentrates at $0$ (i.e.\
$p=q$), and positive once $p>q$ pushes the differential past $\tau$ with positive probability. \emph{(iii)} At
$q=0,p=1$ the differential is maximal ($b$ deterministically connected to $a$, no distractor connected), the
oracle-clique configuration; at $p=q$ it is $0$ by (i).
\end{proof}

\paragraph{Numerical verification} (\texttt{paperA\_alignment\_sim.py}). $N{=}30$, $q{=}0.05$, Katz kernel,
$B{=}1$: the kernel differential rises $-0.001\to0.113$ and the graph advantage over the embedding kernel rises
$+0.007\to+0.115$ as $p-q:0\to0.95$---monotone, essentially $0$ at $p=q$---while the embedding kernel stays at
chance ($0.500$) throughout, unable to surface the buried bridge. The advantage is linear in $p-q$ (slope
$\approx0.11$), matching Theorem~\ref{thm:align}(i)--(ii).

\begin{remark}[Why UCB, not VOI]
The anchor is judged first \emph{because} UCB weights the prior mean; a myopic value-of-information rule that
scores only expected label information can judge elsewhere and forfeit the propagation that surfaces $b$. This is
the mechanism behind the paper's ``UCB beats VOI'' finding: the value of the anchor judgment is not its own label
but the \emph{downstream} surfacing it enables---an inherently non-myopic quantity.
\end{remark}
\begin{remark}[Measurable alignment]
In practice the alignment $p-q$ is estimated by \emph{gold-connectivity}. Theorem~\ref{thm:align} explains why that
single scalar predicts the structural gain, why the oracle-graph ceiling and the adversarial (MuSiQue) failure are
the two ends of one axis, and why the effect is bounded: the gain saturates as $p-q\to$ its clique maximum.
\end{remark}
\end{document}
